\begin{frame}{NewSQL : le meilleur des deux mondes ?}

    NewSQL est un terme générique pour désigner les \textbf{bases de données relationnelles distribuées} qui combinent les \textbf{avantages des bases de données relationnelles} (ACID, SQL) avec ceux des \textbf{bases de données NoSQL} (scalabilité, haute disponibilité).

    \vspace{0.8cm}

    \underline{Exemples} : Google Spanner, CockroachDB, NuoDB, VoltDB, MemSQL, Clustrix, etc.

    \note{
        NewSQL répond au besoin de conserver un modèle relationnel et des garanties
        ACID tout en distribuant le stockage et le calcul. Ce n'est pas une
        combinaison parfaite de tous les avantages : la distribution ajoute
        de la complexité, des coûts et des compromis.

        La question finale à poser est : l'application a-t-elle besoin  de transactions relationnelles fortes ou d'un modèle NoSQL spécialisé ?
    }

\end{frame}
